\section{Why even the minimal affine Fourier buffer is insufficient}
\label{s32:sec:minimal-buffer-obstruction}
The preceding Fourier correction is useful for appropriately weighted,
very accurate inputs. Here we show a limitation that also applies to
the \emph{smallest} nonnegative buffer, rather than just its universal
upper bound. The argument uses elementary low-degree concentration;
it does not exclude other nonlinear correction procedures.

Let $N\ge20$ be a multiple of four, $d=N-4$, $M=2^d$ and
$K=d(d+1)/2$. Let $P$ be uniform on $N$ distinct points of the
$d$-bit cube. For $1\le|S|\le2$, let
$m_S=\mathbb E_P\chi_S$, and put
\[
 f(x)=\sum_{1\le|S|\le2}m_S\chi_S(x),\qquad
 \sigma^2=\sum_Sm_S^2.
\]
For any feasible buffer $\beta\ge0$, consider
\[
 Q_\beta(x)=\frac{P(x)+(\beta-f(x))/M}{1+\beta}.
\]
Its least nonnegative choice is
$\beta_* =\max(0,\max_x(f(x)-MP(x)))$.

\begin{theorem}[Sparse uniform candidates cannot enter the extraction window]
\label{s32:thm:minimal-buffer-obstruction}
If $P$ is not already unbiased and pairwise independent, then every
nonnegative correction $Q_\beta$ in this family satisfies
\[
 \frac1N-C(Q_\beta)\ge\frac1{8N^2\sqrt K}>
 \frac1{8N^3}>\frac1{1024N^4}.
\]
In particular none enters the sufficient window proved in round six.
If $P$ is already exactly isotropic, its uniform $N$ atoms already
give a Hadamard matrix by the residue equality theorem.
\end{theorem}
\begin{proof}
All expectations in the next paragraph are under the uniform cube.
Fourier orthogonality gives $\mathbb E f=0$ and
$\mathbb E f^2=\sigma^2$. Cauchy--Schwarz gives
$\|f\|_\infty\le\sqrt K\sigma$, so
\[
 \mathbb E|f|\ge\frac{\sigma}{\sqrt K},\qquad
 \mathbb E f_+=\tfrac12\mathbb E|f|
       \ge\frac{\sigma}{2\sqrt K}.
\]
Let $T$ be the $N$-point support of $P$. Its contribution to
$\mathbb E f_+$ is at most $(N/M)\sqrt K\sigma$.
Since $M\ge4NK$, a positive contribution of at least
$\sigma/(4\sqrt K)$ remains outside $T$. Hence
\[
 \beta_*\ge\max_{x\notin T}f(x)
          \ge\frac{\sigma}{4\sqrt K}.
\]
A nonzero $m_S$ is an even integer divided by $N$; therefore
$\sigma\ge2/N$, and
$\beta\ge\beta_0:=1/(2N\sqrt K)$.

The exact collision identity for the correction, valid for every
feasible $\beta$, is
\[
 C(Q_\beta)=\frac1M+
 \frac{1/N-(1+\sigma^2)/M}{(1+\beta)^2}
 \le\frac1M+\frac{1/N-1/M}{(1+\beta_0)^2}.
\]
Here the negative $\sigma^2$ term was discarded \emph{before}
comparing denominators, so no unproved sign assumption is needed.
Since $M\ge2N$, $\beta_0\le1$, and
$1-(1+\beta_0)^{-2}\ge\beta_0/2$, subtraction from $1/N$ gives
the claimed lower bound. Also $\sqrt K<N$.

For completeness, $2^{16}=65536\ge4\cdot20\cdot136$ verifies
$M\ge4NK$ at $N=20$. Increasing $N$ by four multiplies $M$ by
$16$, while $NK$ grows by a factor at most
$(24/20)(20\cdot21)/(16\cdot17)=63/34<2$.
Thus the inequality holds for all stated $N$.
\end{proof}

The conclusion concerns sparse \emph{uniform} starting laws. It does
not negate the correction theorem for nonuniform weighted candidates
whose initial concentration and moment accuracy are both sufficient.
Finding such candidates is still an existence problem, not a
consequence of the correction identity.

\paragraph{Formalization.}
The full theorem is formalized for $d\ge16$ with $N=d+4$ a multiple
of four. The independent specification includes both the zero-defect
Hadamard conclusion and the universal collision obstruction for every
feasible buffer. In particular the proof does not assume the buffer
dominates the Fourier polynomial on the original support. It derives
the necessary off-support bound from the corrected law's actual
nonnegativity, counts the low characters exactly, proves the cube-size
estimate, and uses exact sign-moment quantization. The companion
\texttt{BufferChallenge} has no definition holes. The new proof modules
are compiled with a one-gigabyte per-module memory limit.
